%==============================================================================
% Bernard Benson, PhD --- Academic Curriculum Vitae (long form)
% Compile with: tectonic CV_BernardBenson_Academic.tex
%           or: latexmk -pdf CV_BernardBenson_Academic.tex
%
% Items marked with \verify{...} were inferred from the existing CV and need
% confirmation. They print in red. When all are confirmed, change the
% definition of \verify below to:  \newcommand{\verify}[1]{#1}
%==============================================================================
\documentclass[11pt,letterpaper]{article}

\usepackage[T1]{fontenc}
\usepackage[utf8]{inputenc}
\usepackage{lmodern}
\usepackage[margin=1in]{geometry}
\usepackage{titlesec}
\usepackage{enumitem}
\usepackage{xcolor}
\usepackage{tabularx}
\usepackage{fancyhdr}
\usepackage{lastpage}
\usepackage[hidelinks]{hyperref}

\hypersetup{colorlinks=true, urlcolor=black, linkcolor=black,
  pdftitle={Curriculum Vitae --- Bernard Benson}, pdfauthor={Bernard Benson}}

% --- Unverified content marker (red until confirmed) ------------------------
\newcommand{\verify}[1]{\textcolor{red}{#1}}

% --- Section heading: small-caps with a full-width rule underneath ----------
\titleformat{\section}
  {\large\scshape\bfseries}
  {}{0em}{}[\vspace{-0.5ex}\titlerule]
\titlespacing{\section}{0pt}{1.6ex plus 0.4ex minus 0.2ex}{1.0ex}

% --- Subsection heading: bold italic, no rule -------------------------------
\titleformat{\subsection}
  {\normalsize\bfseries\itshape}
  {}{0em}{}
\titlespacing{\subsection}{0pt}{1.0ex plus 0.2ex}{0.4ex}

\setlist[itemize]{leftmargin=1.4em, topsep=0.2ex, itemsep=0.2ex, parsep=0pt}

\setlength{\parindent}{0pt}

% --- Page footer: name and page X of Y (standard for multi-page CVs) --------
\pagestyle{fancy}
\fancyhf{}
\renewcommand{\headrulewidth}{0pt}
\fancyfoot[L]{\footnotesize Bernard Benson --- Curriculum Vitae}
\fancyfoot[R]{\footnotesize Page \thepage\ of \pageref{LastPage}}
\fancypagestyle{plain}{\fancyhf{}\renewcommand{\headrulewidth}{0pt}%
  \fancyfoot[R]{\footnotesize Page \thepage\ of \pageref{LastPage}}}

% --- Helper: an experience/appointment entry --------------------------------
% \cventry{Role}{Organization, Location}{Dates}
\newcommand{\cventry}[3]{%
  \par\vspace{0.6ex}%
  \noindent\begin{tabularx}{\textwidth}{@{}X r@{}}
    \textbf{#1} & \textbf{#3}\\
    \textit{#2} & \\
  \end{tabularx}%
  \vspace{-0.4ex}%
}

% --- Helper: a simple two-column line (left text, right-aligned date) --------
\newcommand{\daterow}[2]{%
  \noindent\begin{tabularx}{\textwidth}{@{}X r@{}}#1 & #2\end{tabularx}\par
}

% --- Helper: a publication item; bolds the author name automatically ---------
\newcommand{\me}{\textbf{B.~Benson}}
\newlist{pubs}{enumerate}{1}
\setlist[pubs]{leftmargin=2.4em, label=[\arabic*], topsep=0.4ex, itemsep=0.6ex, parsep=0pt, labelsep=0.5em}

%==============================================================================
\begin{document}
\thispagestyle{plain}

% ----------------------------- HEADER ---------------------------------------
\begin{center}
  {\LARGE\bfseries Bernard Benson, PhD}\\[0.5ex]
  {\normalsize Research Scientist, NASA Office of Data Science and Informatics (ODSI)\\
   University of Alabama in Huntsville, Huntsville, AL}\\[0.6ex]
  {\small
    \href{mailto:bernardvbenson@gmail.com}{bernardvbenson@gmail.com}
    \;$|$\; 803-931-2777
    \;$|$\; \href{https://www.linkedin.com/in/bernardbenson}{LinkedIn}
    \;$|$\; \href{https://scholar.google.com/citations?user=uiHWfLEAAAAJ}{Google Scholar}
    \;$|$\; \href{https://github.com/bernardbenson}{GitHub}
  }
\end{center}
\vspace{0.4ex}

% ------------------------- RESEARCH INTERESTS -------------------------------
\section{Research Interests}
My research develops machine-learning methods for scientific data, with a focus on
heliophysics and space weather. I have applied deep neural networks to coronal magnetic-field
modeling, the estimation of magnetic-field parameters from solar images, solar-cycle
prediction, solar-flare forecasting, and coronal mass ejection (CME) arrival-time prediction.
A second line of work uses large language models to make scientific data easier to find and
use: multi-label classification of scientific documents, retrieval-augmented generation (RAG),
and agentic, Model Context Protocol (MCP)-based tools for NASA data archives.

\medskip
\textbf{Keywords:} machine learning for the physical sciences; heliophysics and space weather;
time-series forecasting; uncertainty quantification and ensemble methods; representation
learning; natural language processing and large language models; scientific information
retrieval; computer vision.

% ------------------------------ EDUCATION -----------------------------------
\section{Education}
\cventry{Ph.D., Electrical and Computer Engineering}{University of Alabama in Huntsville, Huntsville, AL}{}

\cventry{M.S., Electrical and Computer Engineering}{University of Alabama in Huntsville, Huntsville, AL}{}

\cventry{B.Tech., Electrical and Electronics Engineering}{Geethanjali College of Engineering and Technology, India}{}

% ----------------------- RESEARCH APPOINTMENTS ------------------------------
\section{Research Appointments}

\cventry{Research Scientist VI}{NASA Office of Data Science and Informatics (ODSI), University of Alabama in Huntsville, Huntsville, AL}{May 2025 -- Present}
\begin{itemize}
  \item Research and develop ensemble large language models for multi-label classification of scientific documents, to support the curation of NASA's Science Discovery Engine (SDE).
  \item Designed a search layer on OpenSearch that modernizes the SDE backend for discovery across NASA science data and literature.
  \item Develop Model Context Protocol (MCP) tools that let LLM agents search the Planetary Data System (PDS4) archive and the SDE.
\end{itemize}

\cventry{Applied Machine Learning Researcher}{Frontier Development Lab (NASA / SETI Institute), Remote}{Jul 2021 -- Aug 2021}
\begin{itemize}
  \item Member of an interdisciplinary team that modeled the effect of thermospheric (solar) drag on satellite orbits, in a NASA-sponsored research sprint.
  \item Developed deep-neural-network ensembles for simultaneous multivariate forecasts of space-weather indices; the work resulted in three NeurIPS 2021 workshop papers.
  \item The team received the Frontier Development Lab Team Excellence Award.
\end{itemize}

% ----------------------------- PUBLICATIONS ---------------------------------
\section{Publications}
{\small Reverse chronological order. The name of the CV holder is in bold.
Full list on \href{https://scholar.google.com/citations?user=uiHWfLEAAAAJ}{Google Scholar}.}

\subsection{Peer-Reviewed Journal Articles}
\begin{pubs}
  \item T.~S.~Newman, C.~W.~Hall, L.~Farris, T.~Singh, N.~Pogorelov, \me, S.~Raza, \textit{et al.}, ``Solar Flare Forecasting Using Machine Learning and SDO/HMI Data: A Multiple Machine Learning Model and Data Curation Technique Comparison Study,'' \textit{The Astrophysical Journal Supplement Series}, vol.~280, no.~2, art.~54, 2025.
  \item T.~Singh, \me, S.~A.~Z.~Raza, T.~K.~Kim, N.~V.~Pogorelov, W.~P.~Smith, and C.~N.~Arge, ``Improving the Arrival-Time Estimates of Coronal Mass Ejections by Using Magnetohydrodynamic Ensemble Modeling, Heliospheric Imager Data, and Machine Learning,'' \textit{The Astrophysical Journal}, vol.~948, no.~2, art.~78, 2023.
  \item \me, W.~D.~Pan, A.~Prasad, G.~A.~Gary, and Q.~Hu, ``On the Estimation of the SHARP Parameter MEANALP from AIA Images Using Deep Neural Networks,'' \textit{Solar Physics}, vol.~296, no.~11, art.~163, 2021.
  \item \me, W.~D.~Pan, A.~Prasad, G.~A.~Gary, and Q.~Hu, ``Forecasting Solar Cycle 25 Using Deep Neural Networks,'' \textit{Solar Physics}, vol.~295, no.~5, art.~65, 2020.
  \item \me, W.~D.~Pan, G.~A.~Gary, Q.~Hu, and T.~Staudinger, ``Determining the Parameter for the Linear Force-Free Magnetic Field Model with Multi-Dipolar Configurations Using Deep Neural Networks,'' \textit{Astronomy and Computing}, vol.~26, pp.~50--60, 2019.
\end{pubs}

\subsection{Conference Proceedings and Workshop Papers}
\begin{pubs}[resume]
  \item R.~Hooda, W.~D.~Pan, and \me, ``Transform Decomposition Switching for Efficient Attribute Compression of 3D Point Clouds Using Neural Networks,'' \textit{International Conference on Computational Science and Computational Intelligence (CSCI)}, 2022.
  \item T.~Singh, \me, \textit{et al.}, ``Improving the Arrival-Time Prediction of Coronal Mass Ejections Using Magnetohydrodynamic Ensemble Modeling, Heliospheric Imager Data, and Machine Learning,'' \textit{Proceedings of the 2nd Machine Learning in Heliophysics Conference}, 2022.
  \item \me, E.~Brown, S.~Bonasera, G.~Acciarini, J.~A.~P\'erez-Hern\'andez, \textit{et al.}, ``Simultaneous Multivariate Forecast of Space Weather Indices Using Deep Neural Network Ensembles,'' \textit{Fourth Workshop on Machine Learning and the Physical Sciences, NeurIPS 2021}, 2021.
  \item S.~Bonasera, G.~Acciarini, J.~A.~P\'erez-Hern\'andez, \me, E.~Brown, \textit{et al.}, ``Dropout and Ensemble Networks for Thermospheric Density Uncertainty Estimation,'' \textit{Bayesian Deep Learning Workshop, NeurIPS 2021}, 2021.
  \item E.~J.~E.~Brown, S.~Bonasera, \me, J.~A.~P\'erez-Hern\'andez, G.~Acciarini, \textit{et al.}, ``Learning the Solar Latent Space: Sigma-Variational Autoencoders for Multiple Channel Solar Imaging,'' \textit{Fourth Workshop on Machine Learning and the Physical Sciences, NeurIPS 2021}, 2021.
  \item \me, Z.~Jiang, W.~D.~Pan, G.~A.~Gary, and Q.~Hu, ``Determination of Linear Force-Free Magnetic Field Constant Alpha Using Deep Learning,'' \textit{International Conference on Computational Science and Computational Intelligence (CSCI)}, 2017.
\end{pubs}

\subsection{Conference Abstracts and Presentations}
\begin{pubs}[resume]
  \item C.~Davis, B.~Praveen, E.~Foshee, K.~Bugbee, D.~Sharma, N.~Pantha, S.~Awale, \me, \textit{et al.}, ``LLMs, Encoders, and Agentic Workflows for Curation and Discovery in NASA's Science Discovery Engine,'' \textit{AGU Fall Meeting 2025}, 2025.
  \item T.~Singh, L.~Farris, C.~Hall, T.~Newman, \me, S.~Raza, and N.~Pogorelov, ``Solar Flare Forecasting Using Multiple Machine Learning Models and SDO/HMI Data,'' \textit{SHINE 2024 Workshop}, 2024.
  \item T.~Singh, C.~Hall, T.~Newman, \me, S.~Raza, and N.~Pogorelov, ``Solar Flare Forecasting Using Machine Learning and SDO/HMI Data: A Comparison of Multiple ML Models and AR Parameter Time Series,'' \textit{SHINE 2023 Workshop}, 2023.
\end{pubs}

% ------------------------ SELECTED RESEARCH PROJECTS ------------------------
\section{Research Projects}

\daterow{\textbf{LLMs and Agentic Tools for Scientific Data Discovery}}{\textit{NASA ODSI, 2025 -- Present}}
\begin{itemize}
  \item Ensemble LLM and encoder models for multi-label classification of scientific documents in the NASA Science Discovery Engine; MCP-based tools for agentic search of PDS4 and SDE.
\end{itemize}

\daterow{\textbf{Improving CME Arrival Time with Ensemble Models and Machine Learning}}{\textit{UAH, 2021 -- 2023}}
\begin{itemize}
  \item Combined magnetohydrodynamic ensemble simulations with heliospheric-imager data and regression and neural-network models to improve CME arrival-time predictions.
\end{itemize}

\daterow{\textbf{Simultaneous Multivariate Forecasts of Space-Weather Indices}}{\textit{Frontier Development Lab, 2021}}
\begin{itemize}
  \item LSTM ensembles forecasting multiple space-weather indices over a full solar rotation, with ensemble- and dropout-based uncertainty estimates for thermospheric density.
\end{itemize}

\daterow{\textbf{Solar Image Processing for Flare Detection}}{\textit{UAH}}
\begin{itemize}
  \item Spatial- and frequency-domain filtering to enhance coronal loops; latent-space visualization with variational autoencoders; analysis of the relation between SHARP parameters and flaring events.
\end{itemize}

\daterow{\textbf{Time-Series Forecast of Solar Cycle 25 Using Deep Neural Networks}}{\textit{UAH, 2019 -- 2020}}
\begin{itemize}
  \item WaveNet and LSTM models predicting the amplitude and timing of Solar Cycle 25 from historical sunspot records.
\end{itemize}

\daterow{\textbf{Deep Learning for Coronal Magnetic-Field Modeling}}{\textit{UAH, 2017 -- 2021}}
\begin{itemize}
  \item Convolutional neural networks estimating the linear force-free field parameter $\alpha$ for multi-dipolar configurations, and estimating the SHARP parameter MEANALP directly from SDO/AIA images.
\end{itemize}

% ----------------------------- INVITED TALKS --------------------------------
\section{Invited Talks}
\begin{itemize}
  \item ``AI in the Supply Chain,'' College of Business, Mississippi State University, Mar.~2024.
  \item ``Applications of AI in Logistics and Supply Chain,'' Harbert College of Business, Auburn University, Oct.~2023.
  \item ``Deep Neural Networks Applied to Three Problems in Heliophysics: Coronal Magnetic Field Modeling, Free Magnetic Energy, and Forecasting Solar Activity,'' Heliophysics Division, NASA Marshall Space Flight Center, Sep.~2020.
\end{itemize}

% ------------------------------- TEACHING -----------------------------------
\section{Teaching Experience}
\cventry{Graduate Teaching Assistant}{Department of Electrical and Computer Engineering, University of Alabama in Huntsville}{Aug 2016 -- May 2021}
\begin{itemize}
  \item \textbf{CPE 211L --- Introduction to Programming in C++ (Laboratory)}, Lab Instructor, 2016 -- 2021. Led weekly lab sections, taught programming fundamentals.
  \item \textbf{EE 316 --- Electrical Circuits and Electrical Design Laboratory}, Lab Instructor, Summer 2017. Supervised circuit-design experiments and lab safety, and gave technical guidance to students.
\end{itemize}

\medskip
\textbf{Teaching Interests:} machine learning and deep learning; data science for the physical
sciences; natural language processing and large language models; time-series analysis;
signal and image processing; introductory programming (Python, C++).

% --------------------------- HONORS & AWARDS --------------------------------
\section{Honors \& Awards}
\daterow{Team Excellence Award, Frontier Development Lab (NASA / SETI Institute)}{2021}

% --------------------------- PROFESSIONAL SERVICE ---------------------------
\section{Professional Service}
\textbf{Proposal Review}
\begin{itemize}
  \item NASA Heliophysics research proposals.
\end{itemize}
\textbf{Journal and Conference Review}
\begin{itemize}
  \item \textit{Journal of Geophysical Research: Machine Learning and Computation} (AGU).
  \item \textit{Solar Physics} (Springer).
  \item \textit{Journal of Atmospheric and Solar-Terrestrial Physics} (Elsevier).
  \item Journals of the American Astronomical Society.
  \item IEEE SoutheastCon --- AI and ML tracks, Mar.~2021.
\end{itemize}
\textbf{Session Chair}
\begin{itemize}
  \item IEEE SoutheastCon --- Computing track, Mar.~2021.
\end{itemize}

% ------------------------ INDUSTRY EXPERIENCE -------------------------------
\section{Industry Experience}

\cventry{Generative AI Solutions Architect}{Megan Soft VA, Remote}{Dec 2024 -- May 2025}
\begin{itemize}
  \item Developed evaluation methods and metrics for generative-AI chatbots and RAG systems (DeepEval, Ragas), and designed orchestration of multiple LLM agents with overlapping capabilities.
\end{itemize}

\cventry{Data Scientist}{McLeod Software, Birmingham, AL}{Aug 2021 -- Nov 2024}
\begin{itemize}
  \item Led applied ML research for freight-rate prediction across more than 1,100 trucking organizations; built RAG chatbots and LayoutLMv2-based document classifiers.
  \item Gave invited talks at universities and industry conferences on AI in transportation and logistics.
\end{itemize}

\cventry{RF Engineer}{Nexgen Wireless / Global Telecom Associates, Milwaukee, WI}{Mar 2013 -- Mar 2016}
\begin{itemize}
  \item Analyzed drive-test data and network logs to diagnose coverage gaps and interference in cellular networks.
\end{itemize}

% ----------------------------- TECHNICAL SKILLS -----------------------------
\section{Computational Skills}
\begin{itemize}
  \item \textbf{Languages:} Python, C++, SQL, MATLAB.
  \item \textbf{Machine Learning:} PyTorch, TensorFlow, Keras, scikit-learn; CNNs, RNNs/LSTMs, Transformers, VAEs; deep ensembles and uncertainty estimation; LLMs, RAG, and MCP-based agents.
  \item \textbf{Scientific Computing:} NumPy, SciPy, Pandas, Matplotlib, Seaborn, Plotly; solar data products (SDO/HMI SHARP, SDO/AIA).
  \item \textbf{Infrastructure:} OpenSearch, Docker, Git, CI/CD; AWS, Azure, Google Cloud Platform.
\end{itemize}

% ------------------------------- REFERENCES ---------------------------------
\section{References}
Available upon request.

\end{document}
